% CHAPTER 1
\chapter{PENDAHULUAN \& RUANG LINGKUP}

\section{Pendahuluan}

Python adalah bahasa pemrograman yang sangat populer untuk analisis data, machine learning, dan komputasi ilmiah di berbagai bidang. Panduan ini akan fokus pada instalasi Python menggunakan \textbf{Miniconda}, yang merupakan distribusi Python yang ringan namun powerful untuk keperluan scientific computing.
\\
\textbf{Catatan:} Dokumen ini difokuskan \emph{hanya} untuk pengguna macOS.

\section{Persyaratan Sistem}

\subsection{Spesifikasi Minimum}
\begin{itemize}
    \item RAM: 4 GB (disarankan 8 GB atau lebih)
    \item Storage: 5 GB ruang kosong
    \item Koneksi internet untuk download
    \item Sistem operasi yang didukung (fokus dokumen ini):
    \begin{itemize}
        \item macOS 10.15 atau lebih baru (disarankan macOS 12+ pada Apple Silicon/Intel)
    \end{itemize}
\end{itemize}

\chapter{INSTALASI MINICONDA DI MACOS}

\section{Mengapa Miniconda?}

Miniconda adalah distribusi Python yang ideal untuk scientific computing karena:
\begin{itemize}
    \item \textbf{Ringan:} Hanya menginstall komponen essential
    \item \textbf{Fleksibel:} Kontrol penuh terhadap package yang diinstall
    \item \textbf{Environment Management:} Mudah mengelola multiple environment
    \item \textbf{Scientific Focus:} Optimized untuk scientific computing
    \item \textbf{Cross-platform:} Konsisten lintas OS; panduan ini memfokuskan langkah pada macOS
\end{itemize}

\section{Instalasi Miniconda di macOS}

% Only macOS install steps remain, Windows and Linux blocks removed
% \subsubsection{macOS} % removed as per instruction
\begin{enumerate}
    \item Download installer dari \url{https://docs.conda.io/en/latest/miniconda.html} (Lihat Gambar~\ref{fig:macos-download})
    \item Pilih installer sesuai arsitektur: \textbf{arm64 \(Apple Silicon\)} atau \textbf{x86\_64 \(Intel\)}.
    \item Install menggunakan salah satu metode:
    
    \textbf{Metode 1: Graphical Installer}
    \begin{itemize}
        \item Download file \texttt{.pkg}
        \item Double-click dan ikuti wizard instalasi
    \end{itemize}
    
    \textbf{Metode 2: Command Line (Recommended)}
    \begin{lstlisting}[language=bash]
# Download installer (pilih salah satu sesuai arsitektur)
# Apple Silicon (arm64)
curl -O https://repo.anaconda.com/miniconda/Miniconda3-latest-MacOSX-arm64.sh
# Intel (x86_64)
# curl -O https://repo.anaconda.com/miniconda/Miniconda3-latest-MacOSX-x86_64.sh

# Install (ganti nama file sesuai yang diunduh)
bash Miniconda3-latest-MacOSX-arm64.sh

# Saat ditanya "Do you wish the installer to initialize Miniconda3", jawab "yes"
    \end{lstlisting}
Lihat Gambar~\ref{fig:macos-terminal-curl} untuk contoh tampilan Terminal saat mengunduh dan menjalankan installer.
    
    \item Restart terminal atau jalankan:
    \begin{lstlisting}[language=bash]
source ~/.zshrc  # untuk zsh
# atau
source ~/.bash_profile  # untuk bash
conda init zsh  # atau bash/fish sesuai shell
    \end{lstlisting}
Lihat Gambar~\ref{fig:macos-conda-init}.
    
    \item Verifikasi instalasi:
    \begin{lstlisting}[language=bash]
conda --version
python --version
which python
    \end{lstlisting}
Contoh verifikasi pada Gambar~\ref{fig:macos-conda-version}.
\end{enumerate}

% ==== Placeholder Gambar Proses Instalasi macOS ====
% ==== Placeholder Gambar Proses Instalasi macOS ====
\begin{figure}[t]
  \centering
  \includegraphics[width=0.9\linewidth]{Images/ss-20.05.37.png}
  \caption{Contoh halaman unduhan Miniconda untuk macOS dengan pilihan arsitektur (arm64 untuk Apple Silicon atau x86\_64 untuk Intel).}
  \label{fig:macos-download}
\end{figure}

\begin{figure}[t]
  \centering
  \includegraphics[width=0.9\linewidth]{Images/ss-20.09.44.png}
  \caption{Contoh tampilan Terminal ketika menjalankan perintah \texttt{curl} dan \texttt{bash} untuk memasang Miniconda.}
  \label{fig:macos-terminal-curl}
\end{figure}

\begin{figure}[t]
  \centering
  \includegraphics[width=0.9\linewidth]{Images/ss-20.11.40.png}
  \caption{Contoh proses inisialisasi shell menggunakan \texttt{conda init} dan pemuatan \texttt{~/.zshrc}.}
  \label{fig:macos-conda-init}
\end{figure}

\begin{figure}[t]
  \centering
  \includegraphics[width=0.9\linewidth]{Images/ss-20.14.21.png}
  \caption{Verifikasi \texttt{conda --version}, \texttt{python --version}, dan \texttt{which python}.}
  \label{fig:macos-conda-version}
\end{figure}

% ==== END Placeholder Gambar Proses Instalasi macOS ====


\chapter{KONFIGURASI \& LINGKUNGAN PENGEMBANGAN}

\section{Konfigurasi Awal Miniconda}

Setelah Miniconda terinstall, kita akan setup environment untuk belajar (untuk tugas ini, hanya paket NumPy, SymPy, dan SciPy yang perlu diinstal):

\begin{lstlisting}[language=bash]
# Update conda ke versi terbaru
conda update -n base -c defaults conda

# Buat environment baru untuk belajar
conda create --name belajar python=3.13.0

# Aktivasi environment
conda activate belajar

# Rekomendasi: aktifkan channel conda-forge dan strict priority
conda config --add channels conda-forge
conda config --set channel_priority strict

# Install paket yang ditugaskan
conda install python=3.13.0 numpy sympy scipy

# Verifikasi instalasi
python -c "import numpy, sympy, scipy; print('Assigned packages OK')"
\end{lstlisting}
\begin{figure}[t]
  \centering
  \includegraphics[width=0.9\linewidth]{Images/ss-20.14.21.png}
  \caption{Pembuatan environment \texttt{belajar} dan instalasi paket inti.}
  \label{fig:macos-env-setup}
\end{figure}


\subsection{Mengelola Environment}

\begin{lstlisting}[language=bash]
# List semua environment
conda env list

# Aktivasi environment
conda activate belajar

# Deaktivasi environment
conda deactivate

# Export environment untuk sharing
conda env export > environment.yml
# atau versi ringkas (tanpa build hash): conda env export --from-history > environment.yml

# Create environment dari file
conda env create -f environment.yml

# Remove environment
conda env remove --name belajar
\end{lstlisting}

\section{Setup Development Environment}

\subsection{Jupyter Notebook untuk Analisis Interaktif}

Jupyter Notebook sangat penting untuk pembelajaran dan eksperimen interaktif:

\begin{lstlisting}[language=bash]
# Install Jupyter di environment belajar
conda activate belajar
conda install jupyterlab notebook

# Tambahkan kernel ke Jupyter
python -m ipykernel install --user --name belajar --display-name "Python (Belajar)"

# Jalankan Jupyter Notebook
jupyter notebook

# Atau JupyterLab (interface modern)
jupyter lab
\end{lstlisting}
\begin{figure}[t]
  \centering
  %\includegraphics[width=0.9\linewidth]{images/macos-06-jupyterlab.png}
  \caption{JupyterLab berjalan di peramban dengan kernel \texttt{Python (Belajar)}.}
  \label{fig:macos-jlab}
\end{figure}

\subsection{IDE yang Direkomendasikan}

\subsubsection{VS Code dengan Extension}
\begin{enumerate}
    \item Download VS Code dari \url{https://code.visualstudio.com/}
    \item Install extension berikut:
    \begin{itemize}
        \item Python (Microsoft)
        \item Jupyter (Microsoft)
        \item Python Docstring Generator
        \item Pylint
        \item GitLens (untuk version control)
    \end{itemize}
    \item Configure Python interpreter:
    \begin{itemize}
        \item Tekan \texttt{Ctrl+Shift+P} (Windows/Linux) atau \texttt{Cmd+Shift+P} (macOS)
        \item Pilih "Python: Select Interpreter"
        \item Pilih interpreter dari environment belajar
    \end{itemize}
\end{enumerate}

\subsubsection{Spyder (Scientific Python IDE)}
\begin{lstlisting}[language=bash]
conda activate belajar
conda install spyder
spyder
\end{lstlisting}

\section{Package Management Best Practices}

\subsubsection{Menggunakan environment.yml}

Buat file \texttt{environment.yml} minimal sebagai berikut:

\begin{lstlisting}
dependencies:
  - python=3.13.0
  - numpy
  - sympy
  - scipy
  - jupyter
  - ipykernel
\end{lstlisting}

\subsubsection{Install dari environment.yml}
\begin{lstlisting}[language=bash]
conda env create -f environment.yml
conda activate geofisika
\end{lstlisting}

% ===== Inserted: INSTAL SYMPY & SCIPY, VS CODE TUTORIAL =====

\chapter{APLIKASI, TROUBLESHOOTING, DAN PENUTUP}

\section{Instal SymPy dan SciPy}
\subsection{Menggunakan pip}
\begin{lstlisting}[language=bash]
# Pastikan berada di environment "belajar"
conda activate belajar

# Perbarui pip (opsional)
python -m pip install --upgrade pip

# Instal SymPy dan SciPy via pip
pip install sympy scipy

# Verifikasi versi
python -c "import sympy, scipy; print(sympy.__version__, scipy.__version__)"
\end{lstlisting}

\subsection{Menggunakan conda}
\begin{lstlisting}[language=bash]
# Pastikan berada di environment "belajar"
conda activate belajar

# Instal SymPy dan SciPy via conda
conda install sympy scipy

# Verifikasi versi
python -c "import sympy, scipy; print(sympy.__version__, scipy.__version__)"
\end{lstlisting}

\section{Menjalankan Python di VS Code}
\begin{enumerate}
  \item Instal VS Code dari \url{https://code.visualstudio.com/} dan ekstensi \textbf{Python} (Microsoft).
  \item Buka folder kerja baru, buat berkas \texttt{hello.py} berisi:
\begin{lstlisting}[language=Python]
print("Belajar program python")
\end{lstlisting}
  \item Pilih interpreter: tekan \texttt{Cmd+Shift+P} \(macOS\) → ketik "Python: Select Interpreter" → pilih interpreter dari environment \texttt{belajar}.
  \item Jalankan: klik tombol \textbf{Run} di kanan atas editor, atau tekan \texttt{Ctrl+F5}.
  \item Output "Belajar program python" akan muncul di panel Terminal/Output VS Code.
\end{enumerate}

\begin{figure}[t]
\centering
\includegraphics[width=0.9\linewidth]{Images/ss-20.17.11.png}
\caption{Tampilan VS Code saat menjalankan program \texttt{belajar.py} di environment \texttt{belajar}.}
\label{fig:vscode-run}
\end{figure}

% ====== END Inserted Section =======

\subsection{Test Script Python}

Buat file \texttt{test\_installation.py} dengan konten berikut:

\begin{lstlisting}[language=Python]
#!/usr/bin/env python3
"""
Script untuk test instalasi Python dan package essentials
"""

import sys
import platform

def test_python_installation():
    print("=== Python Installation Test ===")
    print(f"Python Version: {sys.version}")
    print(f"Platform: {platform.platform()}")
    print(f"Architecture: {platform.architecture()}")
    
def test_essential_packages():
    print("\n=== Testing Essential Packages ===")
    packages = ['numpy', 'sympy', 'scipy']
    
    for package in packages:
        try:
            __import__(package)
            print(f"[OK] {package} - OK")
        except ImportError:
            print(f"[ERROR] {package} - NOT FOUND")

if __name__ == "__main__":
    test_python_installation()
    test_essential_packages()
    print("\n=== Test Completed ===")
\end{lstlisting}

Jalankan script dengan:
\begin{lstlisting}[language=bash]
python test_installation.py
\end{lstlisting}

\section{Troubleshooting Common Issues}

% Remove Windows Issues and Linux Issues blocks (macOS Issues retained)
\subsection{macOS Issues}
\begin{itemize}
    \item \textbf{Certificate errors:} Untuk Conda, sertifikat biasanya sudah terkelola. Jika ada error SSL, perbarui paket CA: \texttt{conda install certifi ca-certificates}
    \item \textbf{Command not found:} Gunakan \texttt{python3} dan \texttt{pip3}
\end{itemize}

\section{Best Practices}

\subsection{Environment Management}
\begin{enumerate}
    \item Selalu gunakan virtual environment untuk setiap project
    \item Dokumentasikan dependencies dalam \texttt{environment.yml} (preferable untuk conda)
    \item Gunakan version control (Git) untuk project code
    \item Regular update package menggunakan \texttt{pip list --outdated}
\end{enumerate}

\subsection{Security Considerations}
\begin{enumerate}
    \item Download Python hanya dari sumber official
    \item Verifikasi checksum file installer
    \item Hati-hati dengan package dari PyPI yang tidak terverifikasi
    \item Gunakan \texttt{pip-audit} untuk security scanning
    \item Gunakan Apple Silicon arm64 build jika Mac Anda M1/M2/M3; hindari mencampur paket x86\_64 melalui Rosetta kecuali sangat diperlukan.
\end{enumerate}

\section{Resources dan Referensi}

\subsection{Official Documentation}
\begin{itemize}
    \item Python Official: \url{https://docs.python.org/}
    \item pip Documentation: \url{https://pip.pypa.io/}
    \item Virtual Environments: \url{https://docs.python.org/3/tutorial/venv.html}
\end{itemize}

\subsection{Learning Resources}
\begin{itemize}
    \item Python.org Tutorial: \url{https://docs.python.org/3/tutorial/}
    \item Real Python: \url{https://realpython.com/}
    \item Python Package Index: \url{https://pypi.org/}
\end{itemize}

\subsection{Community}
\begin{itemize}
    \item Python Community: \url{https://www.python.org/community/}
    \item Stack Overflow: \url{https://stackoverflow.com/questions/tagged/python}
    \item Reddit: \url{https://www.reddit.com/r/Python/}
\end{itemize}

\section{Kesimpulan}

Instalasi Python dengan \textbf{Miniconda} merupakan pilihan optimal karena:

\begin{itemize}
    \item \textbf{Lightweight dan Efisien:} Hanya menginstall komponen yang diperlukan.
    \item \textbf{Environment Management:} Mudah mengelola multiple project dengan dependencies berbeda.
    \item \textbf{Scientific Computing:} Dioptimalkan untuk paket sains umum seperti NumPy dan SciPy.
\end{itemize}

\subsection{Workflow Recommended}

\begin{enumerate}
    \item Install Miniconda
    \item Buat environment khusus untuk setiap project: \texttt{conda create --name project\_name python=3.13.0}
    \item Install package yang diperlukan: \texttt{conda install package\_name}
    \item Dokumentasikan environment: \texttt{conda env export > environment.yml}
    \item Gunakan Jupyter Notebook untuk analisis interaktif
    \item Gunakan VS Code atau Spyder untuk development script
\end{enumerate}

\subsection{Next Steps}

Setelah setup selesai, Anda dapat:
\begin{itemize}
    \item Lanjutkan latihan dasar Python dan NumPy/SciPy.
    \item Gunakan Jupyter Notebook untuk eksplorasi interaktif.
    \item Kembangkan proyek kecil dengan struktur environment yang terdokumentasi.
\end{itemize}
